% Exact Stavi operators and the completed positive partition foundation.

\section{Proper gaps and the Stavi partition constructions}
\label{sec:annex-p09-stavi-gap-temporal-logic}

\begin{annexitem}{annex:p9.stavi.gap.cuts}{thm:p09-stavi-gap-cuts}{The exact cut characterization}
\begin{proof}
A downward-closed nonempty proper set \(L\) lies entirely
below its complement \(R\). If \(L\) has no maximum
and \(R\) has no minimum, a proposed supremum \(c\)
of \(L\) cannot lie in either side: in \(L\) it would
be its maximum, and in \(R\) it would be its minimum.
Thus it has no supremum and satisfies the source's
proper-gap definition.

Conversely, a source cut with no supremum has no
maximum. If its complement had a least point \(r\),
then \(r\) would be a supremum of the cut and
would not belong to it, contrary to the cut
condition. Hence it has exactly the two properties
in the body's definition.

In a complete order \(L\) is nonempty and bounded
above by any point of \(R\), so it must have
a supremum; there is therefore no proper gap.
Conversely suppose a nonempty upper-bounded set
\(S\) has no supremum. Its downward closure
\[
                   L=\{t:\exists s\in S,\ t\leq s\}
\]
is nonempty. Any upper bound of \(S\) lies
outside \(L\), since otherwise it would be
a maximum of \(S\) and thus its supremum.
So the complement is nonempty. The set \(S\)
has no maximum; for each \(l\in L\), take
an \(S\)-point at or above \(l\), then a
still larger one. Hence \(L\) has no
maximum. If its complement had a least
point, that point would be the least
upper bound of \(S\), a contradiction.
This constructs a proper gap. Reversing
the order proves the corresponding
infimum assertion. Thus absence of
proper gaps is equivalent to both
Dedekind-completeness clauses.

The assertions about a point next
to a gap follow from absence of
a greatest left or least right
point. For distinct gaps \(g<h\),
choose a point in \(h\setminus g\);
it lies strictly between the
two gaps. Reversal exchanges
the two nonempty sides and
their missing extremal points.
A cut located at an actual
first or last point has an
empty side or a side with
an extremal point, and so
is not proper.
\end{proof}
\end{annexitem}

\begin{annexitem}{annex:p9.stavi.definable.gap}{thm:p09-stavi-definable-gap}{The initial-run gap formula, with existence and uniqueness}
\begin{proof}
Fix a label's truth set \(P\) and a current point \(t\).
Suppose first that the temporal conjunction \(\gamma^+(P)\)
holds. Define
\[
       L=\{u:u\leq t\}\ \cup\
             \{u:t<u,\ \forall v(t<v<u\to P(v))\}.
\]
This is downward closed. The first conjunct \(P\mathbin{\mathsf U_s}P\)
supplies a member of \(L\) strictly above \(t\).
The second conjunct says that every such member
satisfies \(P\); otherwise it would witness
\(P\mathbin{\mathsf U_s}\neg P\).
The third conjunct supplies \(d>t\) with
\(\neg P(d)\), which therefore is not in \(L\).
Both sides of this cut are nonempty.

If \(L\) had a maximum \(c\), it would be
strictly above \(t\), would satisfy \(P(c)\),
and would have \(P\) throughout \((t,c)\).
For every \(r>c\), membership in \(L\) fails
at \(r\), so there is a point \(v\in(t,r)\)
with \(\neg P(v)\). That point must be
above \(c\). Hence \(\partial^+\neg P(c)\)
holds. The point \(c\) would witness
\(P\mathbin{\mathsf U_s}(P\land\partial^+\neg P)\),
contrary to the fourth conjunct. Thus
\(L\) has no maximum.

If the complement \(R\) had a minimum \(d\),
failure of its initial-run condition would
supply \(\neg P(v)\) at some \(t<v<d\).
Every \(L\)-point above \(t\) satisfies \(P\),
so \(v\) would belong to \(R\), contradicting
minimality. Thus the cut is a proper gap
\(g\), with \(t<g\).
The same initial-run argument gives \(P\)
on \((t,g)\). For any \(r\in R\), its
failure to belong to \(L\) supplies
\(\neg P(v)\) at \(t<v<r\), and this
\(v\) belongs to \(R\). Therefore
\(\neg P\) is coinitial to the right
of the gap, as required.

Conversely suppose a \(P\)-gap \(g\)
succeeds \(t\). Choose an actual
\(u\in(t,g)\), using properness.
It gives \(P\mathbin{\mathsf U_s}P\).
A proposed witness for
\(P\mathbin{\mathsf U_s}\neg P\)
cannot be before \(g\), where \(P\)
holds; if it is after \(g\), a
\(\neg P\)-point occurs between
the gap and that witness and
violates its \(P\)-gap requirement.
Thus the second conjunct holds.
The right side is nonempty and
has coinitial failures, giving
\(F\neg P\).

Any point \(u>t\) reachable by
an all-\(P\) open interval must
be before \(g\), by the same
failure argument. At such a
point choose \(v\in(u,g)\).
The whole \((u,v)\) has \(P\),
so \(\partial^+\neg P(u)\)
is false. Consequently no
point witnesses the fourth
conjunct's forbidden until.
All four conditions hold.

The gap is unique. For any
gap satisfying the conditions,
a point on its left belongs
to the displayed initial-run
set \(L\). A point on its
right has a failure strictly
between \(t\) and itself,
so it does not belong to
that set. Thus the left
side of every such gap
is exactly \(L\).

Expanding the two ordinary
until clauses, \(F\), and
the displayed \(\partial^+\)
gives exactly the four
first-order conjuncts in
the body. Every quantified
variable is an actual
domain point. Fresh
renaming makes every
copy capture-free.
No conjunct tests \(P(t)\).
Order reversal gives the
entire past argument,
including uniqueness and
its first-order formula.
\end{proof}
\end{annexitem}

\begin{annexitem}{annex:p9.stavi.first.order.clauses}{thm:p09-stavi-first-order-clauses}{The binary gap truth tables and structural translation}
\begin{proof}
Suppose \(P\mathbin{\mathsf U_{\mathrm{gap}}}Q\)
holds at \(t\), with its unique \(P\)-gap \(g\)
and an actual \(d>g\) such that \(Q\) holds
on \((g,d)\). Coinitial failures of \(P\)
supply \(x_1\in(g,d)\) with \(\neg P(x_1)\).
For any \(y\in(t,x_1)\) with \(\neg P(y)\),
the point \(y\) lies after \(g\), since
\(P\) holds between \(t\) and \(g\).
Every \(z\in(y,x_1)\) lies in \((g,d)\)
and has \(Q\). Thus this \(x_1\) satisfies
the first-order clause.

Conversely suppose that clause holds.
Its first conjunct supplies the \(P\)-gap
\(g\). Its witness \(x_1\), being a
failure of \(P\), lies to the right of
the gap. For an arbitrary \(z\in(g,x_1)\),
choose a failure \(y\in(g,z)\), using
coinitiality. Then \(t<y<z<x_1\);
the universal condition gives \(Q(z)\).
Hence \(Q\) holds throughout \((g,x_1)\).
The point \(d=x_1\) is the required
right bound. This proof does not require
\(Q(x_1)\), nor does it infer a current
point condition from an open interval.

Reverse the order and exchange the sides
of the proper gap. The same argument
gives the past clause. The operand
roles remain unchanged: \(P\) labels
the interval next to the current point,
and \(Q\) labels a stretch beyond the
gap. Reversal exchanges
\(\mathsf U_s,\mathsf S_s\) and
\(\mathsf U_{\mathrm{gap}},\mathsf S_{\mathrm{gap}}\),
commutes with Boolean operations and
fixes proposition letters. Induction
on formulas proves the satisfaction
identity with the opposite model.

The ordinary strict modalities already
have their displayed first-order
witness/gap clauses. For a gap
modality use the formula just proved
and substitute the first-order images
of its two operands, with every
bound variable fresh. Atoms and
Boolean clauses commute with
translation. Structural induction
therefore proves the finite
first-order translation for the
entire four-operator language,
on all nonempty linear orders.
The apparent semantic quantification
over a gap has been eliminated.

A complete order has no proper gap,
so both gap modalities are false
there. No proper gap can lie
after a last point or before a
first point. This proves the
stated endpoint failures.
For the rational example, the
irrational cut has neither a
greatest rational below it nor
a least rational above it.
Both side labels have the
required values, and every
right point has an earlier
right failure of \(P\).
The same cut continues to
satisfy the formula if only
\(P(0)\) is changed.

For \(\mathbb Z+\mathbb Z\),
the two copies form a proper
gap because the first has no
maximum and the second no
minimum. Each point has its
usual neighbors inside its own
copy, so the order is discrete.
At any point of the first
copy, \(P\) holds up to the
gap, failures are coinitial
in the second copy and \(Q\)
holds there. Thus the future
gap formula is true. This
is why discreteness must
not be substituted for
Dedekind completeness in
the vanishing claim.
\end{proof}
\end{annexitem}

\begin{annexitem}{annex:p9.stavi.two.block}{thm:p09-stavi-two-block}{Every boundary type of a two-cell partition}
\begin{proof}
We first verify that each displayed alternative
constructs the required partition.
For \(P\mathbin{\mathsf U_s}Q\) choose
its witness \(b>t\); the cells are
\([t,b)\) and \(\{b\}\), with \(P(t)\)
provided by the outer conjunct.
For \(Q\mathbin{\mathsf U_s}Q\), use
\(\{t\}\) and \((t,b]\).
For
\(P\mathbin{\mathsf U_s}(P\land(Q\mathbin{\mathsf U_s}Q))\),
let \(c>t\) be its first witness and
\(b>c\) its subsequent \(Q\)-witness.
The cells are \([t,c]\) and \((c,b]\).
All are nonempty, and their labels
follow from the strict gap and point
clauses.

For the gap alternative, let \(g>t\)
be the \(P\)-gap and let \(d>g\)
bound its \(Q\)-stretch. Choose
an actual \(b\in(g,d)\), using
the absence of a first point
on the right of a proper gap.
The cells \([t,g)\) and
\((g,b]\) have the required
labels and are nonempty.
This verifies the implication
from the temporal formula.

Conversely suppose \([t,b]\)
has two nonempty cells
\(I_1<I_2\) labelled \(P,Q\).
In particular \(P(t)\).
If \(I_2\) has a first
point \(c\), then \(P\)
holds on \([t,c)\) and
\(Q(c)\), giving the first
until alternative.
If \(I_1\) has a last
point \(c\), then \(I_2=(c,b]\)
has \(Q\) throughout.
When \(c=t\) this gives
the second alternative;
when \(c>t\), \(P\)
on \([t,c]\) gives
the third.

It remains that neither
of these extrema exists.
The cut whose left side
is all points below \(t\)
together with \(I_1\)
is then proper. Call it
\(g\). If any \(u\in I_2\)
has \(P\) throughout
\((t,u)\), that \(Q\)-point
already gives the first
until alternative, possibly
using a different partition.
Otherwise every \(u\in I_2\)
has a failure of \(P\)
strictly between \(t,u\).
Such a failure must lie
in \(I_2\). Hence failures
are coinitial just after
the cut: for a right
point within \(I_2\)
apply that assertion
directly; for a point
after \(b\), use \(b\)
to obtain an earlier
failure. Thus \(g\)
is the \(P\)-gap
succeeding \(t\).
The interval \((g,b)\)
has \(Q\), giving the
Stavi alternative.
All possible boundary
types are covered.
\end{proof}
\end{annexitem}

\begin{annexitem}{annex:p9.stavi.partition.prefix.rays}{thm:p09-stavi-partition-prefix-rays}{The full four-case recurrence and both ray directions}
\begin{proof}
We prove the suffix assertion by backward induction
on \(i\). For the finite seed, \(F_k=D_k\)
means exactly that the singleton interval at the
current point realizes the last-cell suffix.
This is allowed whether or not \(k\) is a flag.
Conversely every point of a realized last cell
has \(D_k\).

Assume the suffix result for \(i\).
In a realizing partition, choose any point
\(t\) of cell \(i-1\). If both adjacent
cells are forced singletons, the next
point \(s\) is an immediate successor
of \(t\), and has \(F_i\) by induction;
thus false-until gives the first case.
If only the first cell is forced
singleton, choose any \(s\) in the
next cell. The interval \((t,s)\)
has \(D_i\), giving the second case.
If only the next cell is forced
singleton, its point \(s\) follows
an all-\(D_{i-1}\) interval from
\(t\), giving the third case.
If neither is forced, truncate the
first cell at \(t\) and the next
cell at any chosen point \(s\)
inside it. The resulting partition
of \([t,s]\) has labels
\(D_{i-1},F_i\), since \(F_i\)
holds throughout the next cell
by induction. The two-block lemma
gives the fourth case.
Every case also has the current
label \(D_{i-1}(t)\).

For the converse, in the first
case the successor witness \(s\)
has its suffix partition by
induction. Prepend \(\{t\}\);
the empty neighbor gap leaves
no uncovered point.
In the second case use
\(\{t\}\) as the first cell,
and extend the first cell of
the suffix partition leftward
by \((t,s)\). That suffix
cell is not forced singleton,
so the extension is permitted,
and every added point has
\(D_i\).
In the third case prepend
\([t,s)\), which has
\(D_{i-1}\), to the suffix
whose first cell is the
forced singleton \(\{s\}\).

In the fourth case, the
two-block lemma supplies
\([t,s]\) partitioned as
\(J_1<J_2\), labelled
\(D_{i-1},F_i\).
Its last point \(s\)
belongs to \(J_2\), so
the induction hypothesis
supplies a suffix partition
starting at \(s\).
Every point of \(J_2\)
satisfies \(F_i\), and
hence \(D_i\).
Merge \(J_2\) with the
first cell of that
suffix at their common
point \(s\). This
cell is not forced
singleton. Prepend
\(J_1\). The result
covers the whole
interval from \(t\)
to the suffix's
finite endpoint,
with all flags
and labels preserved.
This proves both
directions of every
step.

The same induction
works for rays,
with the terminal
seed interpreted
directly. If the
last cell is
unflagged, \(D_k\land GD_k\)
says that the
whole remaining
ray has \(D_k\).
If it is flagged,
\(D_k\land G\bot\)
says that the
current point has
\(D_k\) and no
later point exists,
so the ray is
exactly that
singleton. The
four gluing
operations above
then attach a
ray suffix instead
of a finite one,
with no other
change. This
proves the exact
future-ray formula,
including last
points and every
flag combination.

For the past
version reverse
the order, reverse
the tuple, send
each flag \(i\)
to \(k+1-i\),
and dualize each
input label.
Apply the future
construction and
dualize its output.
The proved modal
duality restores
each label's
original truth
set and exchanges
the two gap
operators.
This is a
finite construction
and proves the
past-ray equivalence
with the same
endpoint care.
\end{proof}
\end{annexitem}

\begin{annexitem}{annex:p9.stavi.positive.partitions}{thm:p09-stavi-positive-partitions}{The finite refinement and quantified-point construction}
\begin{proof}
First consider two partitions of the same nonempty interval,
with cells \(I_1,\ldots,I_k\) and \(J_1,\ldots,J_l\).
Their nonempty intersections are convex and form an ordered
partition of that interval. The sequence of pairs \((i,j)\)
of nonempty intersections starts at \((1,1)\), ends at
\((k,l)\), and advances by \((1,0),(0,1)\) or \((1,1)\).
No row or column can be skipped, since every original cell
is nonempty. There are finitely many such paths.

For each path use the label \(D_i\land E_j\) at
position \((i,j)\). If row \(i\) or column \(j\)
was required singleton, mark that intersection as
singleton and require the corresponding row or
column to occur only once in the path. An actual
singleton cell can meet only one cell of the
other partition. These conditions are therefore
necessary.

Conversely, from a partition realizing one such
path, take the union of its consecutive cells
in each row, and separately in each column.
These are nonempty convex cells and recover
the required labels. A flagged row or column
has just one intersection and that cell is
forced singleton, so its recovered cell has
exactly one point. This proves the intersection
construction in both directions. An unflagged
intersection is allowed to be singleton.
Contradictory label cases may remain as
unsatisfiable disjuncts.

We describe the anchoring operation needed
when endpoints differ or are quantified.
At a named actual point belonging to cell
\(i\), replace that cell by its portion
strictly before the point, the singleton
point itself, and its portion strictly
after it. Enumerate whether each side
portion is empty; retain only nonempty
portions as cells. The middle cell is
forced singleton with label \(D_i\).
If \(i\) was flagged, both side portions
must be empty. Otherwise the retained
side portions have label \(D_i\) and
are unflagged. Every other cell and
its flag are retained with their
new indices.

This operation has finitely many
possibilities, indexed by \(i\)
and the two side-presence choices.
It is exact: a realized partition
determines those choices, and the
reverse operation merges the
pieces back into its original
cell. Keep the positions assigned
to all named points. Two strictly
ordered names cannot occupy the
same forced singleton position.
At a named endpoint of the
interval, the outside portion
is necessarily absent.

In particular, for a closed
interval with distinct actual
endpoints, anchor both as
separate forced singleton
cells. A one-cell unflagged
partition splits into the
two endpoint singletons and
an optional nonempty middle
cell. A one-cell flagged
partition contributes no
case at distinct endpoints.
At coincident endpoints,
only a one-cell partition
can hold, and its condition
is its label at that point.
These facts retain the
strict inequality between
endpoints in later gluing.

Now put a positive Boolean
formula into finite disjunctive
normal form and enumerate
all weak order types of
its named variables. On
each type use one name
for each equality block.
Order and equality atoms
are decided by that type.
A closed partition with
reversed endpoints is
false; a coincident pair
has the just described
one-point condition.
For each remaining
partition field, anchor
every named point in
its interval, taking
all the finite cases.

Split each such anchored
case into fields on
adjacent named closed
intervals and the two
outer rays. Shared
endpoint singletons
are included on both
adjoining fields.
The case retains which
pieces came from each
original cell; gluing
them and merging the
pieces of that cell
recovers the original
partition, with its
singleton requirements.
Conversely any original
partition supplies the
anchored case and those
restricted fields.
Thus the splitting is
an equivalence, not
an independent selection
of incompatible cuts.
Refine every conjunction
of fields on a common
interval by the path
construction above.
This produces finitely
many normal cases with
one partition field on
each adjacent interval
and outer ray. A
whole-order field may
be split at the same
named points; when
there are no names,
refine the whole-order
fields directly.
True neutral fields
have one unflagged
true cell.

To eliminate \(\exists x\),
first use these normal
cases. If \(x\) shares
its equality block
with another surviving
name, remove \(x\)
and keep that value
and all its constraints.
If it is unused,
remove the quantifier,
using nonemptiness
of the domain.
Otherwise its block
has one actual
point between two
neighboring names,
or next to an
outer ray. All
constraints at that
point are already
in its forced
singleton endpoint
cells.

Glue the field ending
at \(x\) and the
field starting there.
Merge their two
singleton copies of
\(x\) into one
forced singleton
labelled by their
conjunction. Retain
all other cells
and flags. The
singleton now has
no free name; its
existence as a
partition cell
chooses the
quantified actual
point. Surviving
endpoints still
occupy their own
distinct forced
singleton cells,
so this point
is strictly between
them whenever the
normal order type
required it.
The same operation
with an outer
ray retains the
one surviving
endpoint, and
with both rays
gives a whole-order
partition when
no name remains.

Conversely a realization
of the glued word
supplies its marked
singleton as \(x\).
Splitting there
recovers the two
original fields
and all labels
at that same
point. It
therefore supplies
exactly the
required existential
witness. Distributing
all finitely many
normal cases proves
closure under
existential quantification.
No negation closure
was used.

Finite conjunction
and disjunction
are included in
these constructions.
Their empty values
can be represented
by whole-order
one-cell partitions
labelled true or
false, respectively;
the order is
nonempty. Every
construction uses
only finite index
lists, Boolean
label conjunctions
and explicit
singleton flags.
This proves
the claimed
effective positive
closure.
\end{proof}
\end{annexitem}

\begin{annexitem}{annex:p9.stavi.basic.partition.translations}{thm:p09-stavi-basic-partition-translations}{First-order definitions and the complete one-variable return}
\begin{proof}
Take a partition expression with labels supplied as
unary predicate names. Its future-ray formula has
the exact semantics already proved. Translate
that formula into first-order logic using the
ordinary and Stavi truth tables. For actual
\(a\leq b\), evaluate it at \(a\) in the
induced chain \([a,b]\): its ray is then
the whole interval \([a,b]\), so it expresses
exactly the bounded partition.

Relativizing every first-order quantifier of
this translation to \([a,b]\) computes that
induced-model value in the original chain.
Include the guard \(a\leq b\); at reversed
endpoints the positive-length partition is
false. If the labels were themselves
temporal formulas, first treat them as
fresh unary names for their whole-model
truth sets. Perform the interval
relativization before substituting
their first-order images. This order
of operations leaves their internal
quantifiers unrestricted, as their
stated semantics requires.

To express an open bounded interval,
pad both excluded actual endpoints:
\[
 \operatorname{Part}(\langle D_1,\ldots,D_k\rangle,O)(a,b)
 =
 \operatorname{Part}(\langle\top,D_1,\ldots,D_k,\top\rangle,
       \{1,k+2\}\cup\{i+1:i\in O\})[a,b].
\]
Each padded endpoint is forced
singleton and has no substantive
label. The one-sided variants
pad only that excluded endpoint.
The exact cell partition is
recovered by removing those
singletons. Empty open intervals
have no positive-length inner
partition and correctly fail.
This proves their first-order
presentations too.

For a simple formula with one
free variable \(z\), its
bounded fields have endpoints
\(z,z\). A one-cell field
is its label at \(z\);
a longer field is false.
Equality and strict-order
atoms on the same name
are true and false.
Future and past ray fields
have the supplied temporal
translations. It remains
to translate a whole-order
partition.

If such a partition exists,
choose any actual \(z\)
in its first nonempty cell.
Its restriction to
\((-\infty,z]\) has a
single cell labelled
\(D_1\), with flag set
\(O\cap\{1\}\).
Its restriction to
\([z,+\infty)\) has the
whole original tuple,
with the first cell
truncated but nonempty.
If the first cell
was forced singleton,
both restrictions have
that singleton at \(z\).
If it was unflagged,
truncation is permitted.

Conversely, suppose those
two partitions exist at
the same point \(z\).
Merge their two first-cell
pieces at \(z\). Their
union is convex and has
\(D_1\), and all other
cells come from the
future partition.
If \(1\in O\), both
pieces are \(\{z\}\),
so the merged cell
is singleton. All
other flags are
unchanged. This
recovers a partition
of the whole order.

The two ray translations
therefore give a
temporal formula \(A\)
such that the whole
partition exists
exactly when \(A\)
holds at some
actual point.
At any current
point this is
\(A\lor FA\lor PA\):
a witnessing point
is equal, later
or earlier, and
each disjunct
supplies such
a witness.
This also covers
a singleton order.
Its first-order
image gives the
whole-order
partition definition.

Substitute these
expressions in
the finite
positive Boolean
formula. This
proves its
temporal return
and, where needed,
its first-order
presentation, without
using any partition
negation theorem.
\end{proof}
\end{annexitem}

\begin{annexitem}{annex:p9.stavi.common.gap}{thm:p09-stavi-common-gap}{Equality of the two initial-run gaps}
\begin{proof}
Suppose both \(\gamma^+(D)\) and
\(\gamma^+(D\land E)\) hold at \(t\).
Let their unique gaps be \(g,h\),
respectively. The initial run of
\(D\land E\) is contained in that
of \(D\), so their left sides
satisfy \(h\subseteq g\).

If \(h<g\), choose an actual
point \(d\) between them.
The label \(D\) holds throughout
\((h,d)\), because this interval
lies between \(t\) and its
\(D\)-gap. Thus the
\((D\land E)\)-gap \(h\)
has a \(D\)-stretch on its
right and witnesses
\((D\land E)\mathbin{\mathsf U_{\mathrm{gap}}}D\).
If instead \(h=g\), failures
of \(D\) are coinitial just
after that gap, so no such
right \(D\)-stretch exists.
Uniqueness of the
\((D\land E)\)-gap prevents
any different gap from
witnessing the formula.
Consequently its negation,
under the two gamma
conditions, expresses
exactly \(h=g\).

At that common gap,
\(D\land E\) holds
throughout \((t,g)\),
which gives the claimed
\(E\)-condition.
Conversely, if a
\(D\)-gap \(g\)
succeeds \(t\)
and \(E\) holds
on \((t,g)\),
then \(D\land E\)
holds there and
every right failure
of \(D\) is also
a failure of
\(D\land E\).
Hence the same
gap is the
\((D\land E)\)-gap,
and the preceding
calculation gives
all three conjuncts
of \(\operatorname{UG}\).
Nothing in the
argument tests
\(E(t)\).
\end{proof}
\end{annexitem}

\begin{annexitem}{annex:p9.stavi.gap.partition}{thm:p09-stavi-gap-partition}{The complete corrected partition cap}
\begin{proof}
First, a partition of \((t,g)\)
into finitely many nonempty cells
cannot have a singleton last cell.
If its last cell were \(\{u\}\),
properness supplies another actual
point \(v\) with \(u<v<g\).
It belongs to none of the
earlier cells, which precede
\(u\), and not to the
last cell. This contradicts
coverage. The same argument
applies to \([t,g)\).
Thus the definition's false
case for a flagged last
cell is exact.

Suppose now \(k\notin O\)
and the required partition
of \((t,g)\) exists,
where \(g\) is the
\(D\)-gap succeeding \(t\).
Every point of that
interval has \(D\).
At every point \(u\)
of its last cell,
\(E_k\) holds throughout
\((u,g)\). The same
gap is the \(D\)-gap
succeeding \(u\),
so the common-gap
lemma gives
\(\operatorname{UG}(D,E_k)(u)\).

Choose any actual
\(b\) in that last
cell. Truncate the
partition at \(b\),
retaining all earlier
cells and the
nonempty initial
portion of the
last cell.
This preserves all
flags because that
last cell is
unflagged. Prepend
the separate
singleton \(\{t\}\)
labelled true.
The resulting
partition of
\([t,b]\)
realizes exactly
the tuple and
shifted flags
in the definition
of \(\operatorname{OG}\).
The finite-prefix
theorem gives
its recursive
formula at \(t\),
and \(\gamma^+(D)\)
already holds.

Conversely suppose
\(\operatorname{OG}\)
holds. Its gamma
conjunct supplies
the unique
\(D\)-gap \(g\).
The finite-prefix
formula supplies
some \([t,b]\)
partition with
the extra initial
singleton and
the \(k\)
nonempty remaining
cells. Thus
\(b>t\).
Every point in
\((t,b]\)
has \(D\),
by the modified
labels.
The point \(b\)
must lie before
\(g\): if it
were after the
gap, a failure
of \(D\) between
the gap and
\(b\) would
contradict that
condition.

The last point
\(b\) has
\(E_k\) and
\(\operatorname{UG}(D,E_k)\).
Its \(D\)-gap
is the same
\(g\), since
\(D\) holds
between \(b\)
and \(g\)
and its
right failures
are unchanged.
Hence \(E_k\)
holds throughout
\((b,g)\).
Extend the
last cell
by that
interval and
remove the
initial singleton
\(\{t\}\).
The result
is the required
partition of
\((t,g)\).
Extending the
last cell is
permitted because
its index is
not in \(O\).
Every earlier
cell and flag
is retained.
This proves
the open case
in both
directions.

For the closed
case, any
partition of
\([t,g)\)
has \(E_1(t)\).
Its first cell
is either the
singleton
\(\{t\}\)
or has further
points. In
the first
case \(k\geq2\),
since a proper
gap has points
strictly above
\(t\) on
its left.
Removing that
cell gives
the second
open alternative
with shifted
flags.
In the
second case
\(1\notin O\).
Removing only
\(t\) leaves
a nonempty
first cell
with the
same label,
so it gives
the first
open alternative.

Conversely, in
the first
open alternative
adjoin \(t\)
to its first
cell. The
union is
convex, its
label follows
from \(E_1(t)\),
and the cell
is unflagged.
In the
second alternative
prepend
\(\{t\}\)
as the
first cell,
which is
allowed whether
or not
index 1
was flagged.
The other
cells and
their correctly
shifted flags
are unchanged.
Both yield
the required
closed partition.

The extra
initial singleton
was labelled
true, not
\(D\).
Thus the
construction
does not
impose \(D(t)\),
which the
open \(D\)-gap
condition does
not require.
The closed
version explicitly
tests \(E_1(t)\),
which its
partition does
require.
These two
point conditions
are distinct.

The constructions
are finite
compositions of
the already
defined modal
formulas and
the finite
prefix recurrence.
Apply order
reversal to
the tuples,
flags and
labels to
obtain the
two past
versions.
The gap
duality proved
earlier preserves
all their
point and
singleton
conditions.
\end{proof}
\end{annexitem}
